% GENERATED FILE. Edit canonical lessons, then run npm run generate:books.
% canonical-source-hash: fnv1a64:394e3e15e4d4ed6d
% canonical-lessons: FR-C404-crevaison, FR-C404-crever, FR-C404-pneu, FR-C404-cric, FR-C404-gonfler

\chapter{Un pneu à plat}
\label{ch:fr-un-pneu-a-plat}
\hlchaptermodality{404}

\begin{chapteropening}
\textbf{By the end of this chapter:} \emph{I can deal with a flat tyre: the puncture, getting one, the tyre, the jack and pumping it up.}

Complete the last lesson of chapter 404: I can deal with a flat tyre: the puncture, getting one, the tyre, the jack and pumping it up.
\end{chapteropening}

\section[la crevaison]{la crevaison — the puncture, the flat tyre}

\label{lesson:FR-C404-crevaison}



\begin{quote}
Before the new one: say the French for the tank, then the French for to start up.
\end{quote}

\subsection*{You'll want to know: la crevaison}
\textbf{la crevaison} — \textquotedblleft{}the puncture, the flat tyre\textquotedblright{}.

\begin{quote}
Nous avons eu une crevaison sur la petite route.
\end{quote}

\begin{quote}
\emph{We had a puncture on the little road.}
\end{quote}

\emph{Réparer une crevaison} — to mend a puncture.

\begin{cousinweb}
\emph{crevaison} is built on \emph{crever} 'to burst' with the ending \emph{-aison} (EtymDB). \emph{Crever} is inherited from Latin \emph{crepāre} 'to crack, to burst'.
\end{cousinweb}

\begin{grammarlens}[title={What you've built}]
One more word for telling, explaining and describing at B1.
\end{grammarlens}

\subsection*{Your turn}
\begin{itemize}
\raggedright
  \item \emph{Say it:} \emph{la crevaison}
  \item \emph{Say it:} \emph{la crevaison}, once more
  \item \emph{Say it:} a sentence of your own with \emph{la crevaison}
  \item \emph{Recall:} say the French for the tank, then the French for to start up, then say \emph{la crevaison} again
\end{itemize}

\begin{tcolorbox}[breakable,colback=teal!4,colframe=teal!35!black,title={Before you move on}]
What does la crevaison mean? (\textquotedblleft{}The puncture, the flat tyre\textquotedblright{}.) Say it once more.
\end{tcolorbox}
\section[crever]{crever — to get a flat tyre; to burst}

\label{lesson:FR-C404-crever}



\begin{quote}
Before the new one: say the French for to start up, then the French for the puncture.
\end{quote}

\subsection*{You'll want to know: crever}
\textbf{crever} — \textquotedblleft{}to get a flat tyre; to burst\textquotedblright{}.

\begin{quote}
J'ai crevé à dix kilomètres de la maison.
\end{quote}

\begin{quote}
\emph{I got a flat tyre ten kilometres from home.}
\end{quote}

Informal: \emph{je suis crevé} — I am worn out.

\begin{cousinweb}
\emph{crever} is inherited through Old French \emph{crever} from Latin \emph{crepāre} 'to crack, to burst' (EtymDB).
\end{cousinweb}

\begin{grammarlens}[title={What you've built}]
One more word for telling, explaining and describing at B1.
\end{grammarlens}

\subsection*{Your turn}
\begin{itemize}
\raggedright
  \item \emph{Say it:} \emph{crever}
  \item \emph{Say it:} \emph{crever}, once more
  \item \emph{Say it:} a sentence of your own with \emph{crever}
  \item \emph{Recall:} say the French for to start up, then the French for the puncture, then say \emph{crever} again
\end{itemize}

\begin{tcolorbox}[breakable,colback=teal!4,colframe=teal!35!black,title={Before you move on}]
What does crever mean? (\textquotedblleft{}To get a flat tyre; to burst\textquotedblright{}.) Say it once more.
\end{tcolorbox}
\section[le pneu]{le pneu — the tyre}

\label{lesson:FR-C404-pneu}



\begin{quote}
Before the new one: say the French for the puncture, then the French for to get a flat tyre.
\end{quote}

\subsection*{You'll want to know: le pneu}
\textbf{le pneu} — \textquotedblleft{}the tyre\textquotedblright{}.

\begin{quote}
Le pneu avant est à plat.
\end{quote}

\begin{quote}
\emph{The front tyre is flat.}
\end{quote}

Plural \emph{les pneus}. \emph{Un pneu neige} — a winter tyre.

\begin{cousinweb}
\emph{pneu} is short for \emph{pneumatique} 'filled with air', a learned word from Latin \emph{pneumaticus}, from Greek \emph{pneûma} 'breath, air' (Lewis \& Short).
\end{cousinweb}

\begin{grammarlens}[title={What you've built}]
One more word for telling, explaining and describing at B1.
\end{grammarlens}

\subsection*{Your turn}
\begin{itemize}
\raggedright
  \item \emph{Say it:} \emph{le pneu}
  \item \emph{Say it:} \emph{le pneu}, once more
  \item \emph{Say it:} a sentence of your own with \emph{le pneu}
  \item \emph{Recall:} say the French for the puncture, then the French for to get a flat tyre, then say \emph{le pneu} again
\end{itemize}

\begin{tcolorbox}[breakable,colback=teal!4,colframe=teal!35!black,title={Before you move on}]
What does le pneu mean? (\textquotedblleft{}The tyre\textquotedblright{}.) Say it once more.
\end{tcolorbox}
\section[le cric]{le cric — the jack (for lifting a car)}

\label{lesson:FR-C404-cric}



\begin{quote}
Before the new one: say the French for to get a flat tyre, then the French for the tyre.
\end{quote}

\subsection*{You'll want to know: le cric}
\textbf{le cric} — \textquotedblleft{}the jack (for lifting a car)\textquotedblright{}.

\begin{quote}
Le cric est sous le tapis, à l'arrière de la voiture.
\end{quote}

\begin{quote}
\emph{The jack is under the mat at the back of the car.}
\end{quote}

The final \emph{c} is pronounced.

\begin{cousinweb}
The origin of \emph{cric} is not settled. It is usually traced to Middle High German \emph{kriec} 'winch'; EtymDB does not list it.
\end{cousinweb}

\begin{grammarlens}[title={What you've built}]
One more word for telling, explaining and describing at B1.
\end{grammarlens}

\subsection*{Your turn}
\begin{itemize}
\raggedright
  \item \emph{Say it:} \emph{le cric}
  \item \emph{Say it:} \emph{le cric}, once more
  \item \emph{Say it:} a sentence of your own with \emph{le cric}
  \item \emph{Recall:} say the French for to get a flat tyre, then the French for the tyre, then say \emph{le cric} again
\end{itemize}

\begin{tcolorbox}[breakable,colback=teal!4,colframe=teal!35!black,title={Before you move on}]
What does le cric mean? (\textquotedblleft{}The jack (for lifting a car)\textquotedblright{}.) Say it once more.
\end{tcolorbox}
\section[gonfler]{gonfler — to pump up, to inflate; to swell}

\label{lesson:FR-C404-gonfler}



\begin{quote}
Before the new one: say the French for the tyre, then the French for the jack.
\end{quote}

\subsection*{You'll want to know: gonfler}
\textbf{gonfler} — \textquotedblleft{}to pump up, to inflate; to swell\textquotedblright{}.

\begin{quote}
Au garage, il gonfle les pneus.
\end{quote}

\begin{quote}
\emph{At the garage he pumps up the tyres.}
\end{quote}

Informal: \emph{ça me gonfle} — that gets on my nerves.

\begin{cousinweb}
\emph{gonfler} goes back to Latin \emph{conflāre} 'to blow together', from \emph{flāre} 'to blow' (EtymDB), probably through a southern dialect.
\end{cousinweb}

\begin{grammarlens}[title={What you've built}]
One more word for telling, explaining and describing at B1.
\end{grammarlens}

\subsection*{Your turn}
\begin{itemize}
\raggedright
  \item \emph{Say it:} \emph{gonfler}
  \item \emph{Say it:} \emph{gonfler}, once more
  \item \emph{Say it:} a sentence of your own with \emph{gonfler}
  \item \emph{Recall:} say the French for the tyre, then the French for the jack, then say \emph{gonfler} again
\end{itemize}

\begin{tcolorbox}[breakable,colback=teal!4,colframe=teal!35!black,title={Before you move on}]
What does gonfler mean? (\textquotedblleft{}To pump up, to inflate; to swell\textquotedblright{}.) Say it once more.
\end{tcolorbox}
